% mage-auto-heading: one-document-one-worker
% mage-auto-heading: pages-as-the-unit-of-scale
% mage-auto-heading: a-new-class-of-input
% mage-auto-heading: remove-media-from-the-working-set
% mage-auto-heading: why-one-level-deeper-matters
% mage-auto-heading: memory-did-not-stop-mattering
% mage-auto-heading: why-are-we-still-chunking
% mage-auto-heading: analyze-before-you-build
% mage-auto-heading: dependencies-can-be-redesigned
% mage-auto-heading: the-second-resource-model
% mage-auto-heading: predicting-latency
% mage-auto-heading: predicting-the-cost-of-a-design-change
% mage-auto-heading: use-prediction-error-as-evidence
% mage-auto-heading: what-carries-forward
% mage-auto-heading: models-get-falsified-the-obligation-survives
% mage-auto-heading: the-resulting-model-family
% Authoritative LaTeX manuscript; Markdown is a generated compatibility projection.

% <!-- part-title: MAGE in Motion: Engineering Through Models -->

% <!-- chapter-title: One Problem, Many Models -->

\chapter{One Problem, Many Models}

How do engineers choose a useful model when they do not yet know the architecture they need? Familiar techniques offer several possibilities: reduce the working set, separate structure from payload, represent dependencies, analyze a critical path, bound concurrency, or measure the result. Experience helps an engineer recognize these possibilities; evidence from the system helps decide which one to pursue.

In DocAble, a file that exceeded the system\textquotesingle s operating regime prompted a succession of changes in representation and architecture. The memory problem exposed one relationship, and changing it made another worth examining. Each step made a consequential relation explicit, changed or analyzed the realization, and supplied evidence for what should be modeled next.

Large documents created both memory and execution-time problems. Out-of-memory failure prevented processing altogether, whereas inefficient execution increased latency and cost while still producing a result. Memory therefore came first in this episode. The later service-sizing model recorded that priority in a one-sided safety margin: exhausting memory kills the process, while over-provisioning costs money. Another domain might prioritize safety, latency, or numerical accuracy; the ordering must follow the consequences.

The sequence also runs from outcomes to explanations. "Reduce peak memory" states an objective. It does not explain why memory grows, which parts of the current realization are essential, or which changes should generalize to the next large document.

\section{One Document, One Worker}\label{heading:one-document-one-worker}

DocAble originally treated an entire document as one unit of execution. A long-running Kubernetes worker opened the document, held the state needed for remediation in memory, modified it, and produced the result. \autoref{interlude-execution-unit} traces how that execution unit later changed.

\begin{figure}[htbp]
\centering
\includegraphics[width=\linewidth]{book/latex/assets/interlude-execution-unit-3ac64b0b5d.pdf}
\caption{\emph{Two execution units for document remediation.} (a) Whole-document execution keeps the artifact\textquotesingle s working state and media resident in one long-running worker. (b) Page-chunked execution divides the document across workers, reducing the working state and memory required by any one worker.}
\label{interlude-execution-unit}
\end{figure}

The document therefore served as both the user\textquotesingle s artifact and the worker\textquotesingle s in-memory unit of computation. Media needed by GenAI operations was read from that resident representation and sent to the foundation model as needed.

For the small presentations used to develop the system, this worked well. The architecture was simple, and the document was a natural boundary: a user submitted a document and expected a document back. Nothing in the observed workload justified another unit of execution.

Real use eventually expanded the workload. Teaching presentations carried more slides, images, and other media, and handling the whole document in one process began to put enough pressure on worker memory that the execution unit needed to shrink.

\section{Pages as the Unit of Scale}\label{heading:pages-as-the-unit-of-scale}

DocAble\textquotesingle s response was ChunkWorker, which divided a document into page ranges (\autoref{interlude-execution-unit}, panel b).

The new representation treated contiguous pages as independent work partitions. For a presentation, a chunk contained a range of slides together with the document state needed to remediate them; the system later recombined the resulting patches into the complete artifact.

Each worker now handled only part of the document, and multiple chunks could run concurrently. More importantly, chunking reduced the memory required by an individual worker. Memory affects both infrastructure cost and cold-start behavior, so a smaller worker was useful even for documents that did not strictly require chunking to fit.

Chunking worked across a much broader workload: ordinary teaching presentations and hundreds of documents from multiple universities. Page ranges gave each ChunkWorker a practical bound on its working state, and experience provided no reason to replace that design merely because another architecture was possible.

This architecture accommodated the inputs encountered so far. The page boundary gradually became part of how DocAble was understood: a document was divided into chunks because chunks were how large documents were processed. A roughly 250 MB media-heavy presentation eventually exceeded that operating regime.

This architectural tradeoff affects the reasoning horizon described in \magevalue{chapter:problem}. Smaller execution units would reduce memory use, but they could also withhold context needed for a semantic judgment. Chunking did not have to stop at page boundaries: we could continue subdividing the document all the way down to individual words or figures. But finer chunking presented an unattractive tradeoff. As we reduce the amount of material per chunk, we reduce the available semantic context and increase the coordination cost. I asked myself: why did so much of the document have to reside in the worker at all?

\section{A New Class of Input}\label{heading:a-new-class-of-input}

A 250 MB PowerPoint file does not imply that remediation needs 250 MB of document content resident in memory.

Modern Office documents use Office Open XML (OOXML). A \texttt{.pptx}, \texttt{.docx}, or \texttt{.xlsx} file is a package containing many parts. Much of its semantic structure is XML: text, slides, relationships, identifiers, and references among objects. Large binary assets such as images and other media live alongside those XML parts inside the package.

The structural portion is comparatively small. It is largely text, and even an unusually large Office document may require only low megabytes to represent its semantic skeleton. The media payload can be hundreds of megabytes. Most remediation operations need only the structural portion; they may reach into the media, but not every media byte need be resident at once.

Until this point, DocAble had largely accepted the storage representation supplied by the document format as its working representation. For Office documents, that meant opening the OOXML package through the document library and letting its structure and media participate in the in-memory object graph; other formats followed analogous format-native representations. That was reasonable while the input sizes fit the worker. Handling a roughly 250 MB media-heavy file without allocating correspondingly large workers required reconsidering the arrangement itself.

The existing representation coupled the two, and \autoref{interlude-storage-vs-reasoning} contrasts that arrangement with the one that replaced it.

\begin{figure}[htbp]
\centering
\includegraphics[width=\linewidth]{book/latex/assets/interlude-storage-vs-reasoning-65a384adac.pdf}
\caption{\emph{Chunked residency versus skeletonized residency.} (a) Page chunking limits how much document state and associated media a worker holds at once, but media remains part of the worker\textquotesingle s resident representation. (b) Skeletonization keeps only document structure and media references resident. Individual media objects stream from external storage to the consumer on demand, so total media size no longer determines persistent worker memory.}
\label{interlude-storage-vs-reasoning}
\end{figure}

So the architecture carried an implicit scaling relation: total media size drove resident worker state, and resident state drove the risk of running out of memory. Finer chunking could continue reducing how much of the document met this relation at one time, but it could not remove the relation itself---and increasingly fine partitions would introduce coordination costs and erode the semantic context available to remediation. On the observed workload that made no difference; the new file made it decisive.

A residency model made the alternative visible (panel b): the structural pipeline could operate on a lightweight representation while media lived in object storage, and operations that needed media could reach it by reference.

\section{Remove Media from the Working Set}\label{heading:remove-media-from-the-working-set}

The memory problem came from coupling document structure to media residency: total media size increased the worker\textquotesingle s persistent working set even though most remediation operations needed only the document\textquotesingle s structural skeleton. Skeletonization removes that coupling. For OOXML documents, the pipeline streams through the package, moves media into external storage, and leaves stable references in a much smaller structural skeleton.

Remediation then operates on the skeleton. If a GenAI operation needs an image, the worker follows the reference and streams that media through the network to the GenAI service without first materializing the document\textquotesingle s complete media payload in worker memory. Recombination follows the same discipline in reverse. The worker no longer needs memory proportional to total media size.

The corresponding memory invariant states the property directly:

% mage-code-language: text
\begin{verbatim}
peak worker RSS <= C
\end{verbatim}

where, within the declared operating envelope,

% mage-code-language: text
\begin{verbatim}
C =
    rolling stream buffer
  + structural skeleton
  + handles
  + remediation working set
\end{verbatim}

and C does not grow with total media size.

Chunking had shown that the observed documents worked. Skeletonization supports a structural claim over a declared class of inputs: peak worker memory is bounded independently of total media size, provided the modeled path never retains the complete payload.

That claim depends on bounded streaming buffers, a structural skeleton whose size remains inside the declared envelope, and bounded active state at each media consumer. The write path must stream too; buffering the complete payload before writing would merely move the memory spike. Pathological structure, oversized individual objects, excessive concurrency, or an unmodeled resource can still exhaust memory.

The implementation enforces the assumptions behind the bound as invariants; measurement checks that the realized system behaves as the model predicts. The PDF path already shows approximately flat peak-RSS scaling across a count-scaled media sweep when each object is flushed; the corresponding OOXML media-size sweep remains a designed validation step.

While this architecture was being developed, agents sometimes proposed removing operations or state that looked expensive but were essential to correct remediation. The fixtures available at the time would not necessarily have exposed every resulting semantic failure. The engineering model supplied a distinction the surface objective did not: which relationships the system depends on, and which resource costs follow from the current realization. Agents could use that structure to correct their account of the system before changing it.

The model also let me judge which relationships an optimization had to preserve while the realization changed faster than I could reconstruct its behavior from code and tests.

Changing the architecture to satisfy the residency model is active alignment. Once correspondence is established, the model and invariant checks provide passive alignment: later changes must preserve the memory property or justify revising the model and its checks. The memory obligation became part of the alignment kernel for the execution changes that followed.

\section{Why One Level Deeper Matters}\label{heading:why-one-level-deeper-matters}

An end-to-end measurement says whether the outcome moved. It does not say why it moved. Peak worker memory is the outcome; media residency is the mechanism that produces it. A model of the mechanism answers the further question: which relationship generated the improvement, and whether that explanation is compatible with the work the system must still perform.

Measurement tests the residency model\textquotesingle s explanation. A peak-memory figure reports what happened on the inputs in one sweep; the model predicts which changes to those inputs should change the result. Comparing the two can expose a missing relationship, as a later sweep did in this case. The engineer needs both: measurements to test the account, and an account of memory growth to assess the next unfamiliar document.

Ousterhout puts the same instruction compactly: measure one level deeper.\citep{ousterhout2018} One level, not arbitrarily deep. The deeper cut earns its cost when it exposes an optimization target, a failure class, or a decision that the surface number hides. Media residency did all three, and the surface number --- peak resident memory --- would have named none of them.

\section{Memory Did Not Stop Mattering}\label{heading:memory-did-not-stop-mattering}

Bringing memory under governance did not make memory disappear. It changed the conditions under which later performance work could proceed. The next questions increasingly concerned latency, throughput, and compute cost, but proposed answers remained subject to the memory relationships already made explicit. A faster design that exceeded the deployment\textquotesingle s memory budget was not an improvement.

That coupling became explicit in the engineering machinery. Concurrency models delegated memory feasibility to the memory model of record rather than doing memory arithmetic of their own. Fan-out width was represented as a peak-memory multiplier, so a throughput author had to size the width against a memory budget. Some throughput changes carried fail-loud guards that refused to run a fan-out whose estimated memory exceeded its budget. A batched-rendering design was chosen over a concurrent one partly because the concurrent alternative measured past the service\textquotesingle s memory ceiling. And in one latency campaign a memory sweep blocked the production deploy altogether. The earlier resource model had become part of the environment within which the later one operated.

\section{Why Are We Still Chunking?}\label{heading:why-are-we-still-chunking}

Page chunking had served two roles at once: it partitioned execution and bounded each worker\textquotesingle s memory. Skeletonization largely removed the memory requirement. That left a separate question: what should determine the unit of execution?

A page tells us where content occurs, not which computations depend on which others. DocAble already represented remediation as a dependency graph. Consider alt-text generation as a concrete example. Each page can first be analyzed independently. Those analyses feed an operation that builds shared document context. Once that context exists, the GenAI queries that generate alt text for individual images become ready to execute independently. The actions and their information dependencies form a remediation graph; the currently executable actions form its frontier.

The same graph admits different execution policies. \autoref{interlude-graph-two-partitions} shows the distinction.

\begin{figure}[htbp]
\centering
\includegraphics[width=\linewidth]{book/latex/assets/interlude-graph-two-partitions-8ab535d9c3.pdf}
\caption{\emph{Two execution policies over the same remediation graph.} The graph contains actions---page analysis, construction of document context, and GenAI queries---and information dependencies among them. (a) A page partition groups actions by the page their artifacts occur on, so it can span several dependency levels. (b) A frontier groups actions by readiness: page analyses execute first, their results enable construction of document context, and that context enables independent GenAI queries for alt text. The work and its dependencies are unchanged; only the execution boundary changes --- so the choice between (a) and (b) is a scheduling decision, not a remediation redesign.}
\label{interlude-graph-two-partitions}
\end{figure}

The graph explains the ordering. Independent work can run together; a consumer waits only for information it actually needs from upstream. Pages do not create those dependencies. Information does. Expressed as a directed acyclic graph, the work currently ready to execute forms a frontier: the nodes with no unmet dependencies. A worker executes that frontier, materializes its outputs, and thereby exposes the next frontier.\magefootnote{topo-dag}{This is ordinary topological traversal of a directed acyclic graph, familiar from compilers, build systems, and dependency schedulers.}

This reframes ChunkWorker. Page-oriented and frontier-oriented execution are not different remediation graphs. They are different ways of partitioning the same graph. Page execution groups work by where it occurs; frontier execution groups work by when it is ready.

Skeletonization removed the memory pressure that had made page ranges the natural execution boundary, so page identity no longer needs to determine execution. It can survive as metadata on graph nodes---useful for locality, debugging, rate limiting, or failure isolation---without remaining the fundamental unit of executable remediation.

For the present problem, the architecture therefore needs only a clean seam between the remediation graph and its executor: the graph supplies ready work and its dependencies; the executor consumes them. The executor can schedule by frontier now and exploit page locality later if measurements justify it. There is no need to build a strategy hierarchy merely because two policies are conceivable. The graph/executor seam preserves the degree of freedom without spending it. In the realized architecture, that execution-driving representation remains distinct from richer analytical projections built over the same computation identities. The distinction matters later: a model may govern realization while another view of the same structure serves analysis.

\section{Analyze Before You Build}\label{heading:analyze-before-you-build}

Representing remediation as a dependency graph does more than provide an execution plan. The graph can be analyzed before an executor --- or a proposed architectural change --- is implemented. \autoref{interlude-analysis-before-build} shows three such analyses.

\begin{figure}[htbp]
\centering
\includegraphics[width=\linewidth]{book/latex/assets/interlude-analysis-before-build-1d6121ea6a.pdf}
\caption{\emph{Analysis before implementation.} (a) In a weighted dependency graph the heaviest chain is the critical path, which sets the latency floor and names an optimization target. (b) Decomposing the dominant node exposes parallel children and can shorten the path, a possibility visible in the model before the split is built. (c) A serial information dependency can be removed by changing the representation supplied to each operation: instead of sibling-to-sibling context propagation, each image consumes one shared precomputed page context, so the alt-text operations run independently.}
\label{interlude-analysis-before-build}
\end{figure}

The weighted graph in panel (a) shows which operations can run concurrently and which must wait. Its heaviest chain from start to finish is the critical path. Adding workers cannot remove the ordering imposed by that chain. If one node dominates the path, the engineer can ask whether its work must remain one operation. Where the semantics permit decomposition, splitting the node may expose parallelism and shorten the path, as panel (b) shows. The model can reveal that candidate before the split is implemented. It cannot establish that the split will pay off: smaller operations may add coordination cost, lose batching efficiency, or contend for another resource. Those costs must also be estimated or measured.

\subsection{Dependencies can be redesigned}\label{heading:dependencies-can-be-redesigned}

The same analysis exposed a problem in contextual image description (panel c). If each image description consumed the description generated for the preceding image, the operations formed a serial chain. For N images, the dependency depth approaches N. If each GenAI call takes roughly t, the serial floor grows as about N·t; adding workers cannot remove that dependency.

The analysis clarified the requirement: each image needed enough surrounding context for an adequate description. It did not need the preceding image\textquotesingle s generated alt text specifically. Rendering the surrounding pages up front supplied a shared visual intermediate representation. Each image could then consume bounded neighboring visual context independently. This increased the payload but removed the serial dependency, allowing the operations to run in parallel. The visual context could also be represented and reused, as textual document context already was. The model thus changed the proposed implementation before it was built.

\subsection{The second resource model}\label{heading:the-second-resource-model}

The critical-path model explained which dependencies could determine latency. A second model captured a different source of cost: repeated whole-document work. Let \emph{R} be the number of remediation rounds, \emph{F} the cost of producing the whole document once, and \emph{G} the remaining GenAI work. The model represented execution without page chunks as

\[
T = (R+1)\,F + G
\]

For example, suppose a document requires four remediation rounds, whole-document production takes 10 seconds, and the remaining GenAI work takes 20 seconds. The model predicts \emph{T} = (4+1)(10) + 20 = 70 seconds. Fifty of those seconds come from producing the document five times: once for each remediation round and once more to apply the accumulated changes. Cutting the GenAI work in half would save 10 seconds; eliminating two production rounds would save 20. The model therefore identifies repeated whole-document production as a quantity worth attacking.

The model was also falsifiable. Its calibration procedure inverted observed execution times to recover its components, and it rejected any fit whose inferred costs came out negative. Measurement could therefore disagree with the engineering account rather than merely supplying another number to optimize against.

\subsection{Predicting latency}\label{heading:predicting-latency}

The same graph supports a first-order latency model. For the representative workload used during this analysis, the relevant GenAI chain had three dependent levels whose per-call times add to an irreducible floor (\autoref{interlude-latency-floor}).

\begin{figure}[htbp]
\centering
\includegraphics[width=\linewidth]{book/latex/assets/interlude-latency-floor-818a3f4133.pdf}
\caption{\emph{The GenAI dependency chain and its latency floor.} A section summary feeds a document summary feeds alt text; because each level waits for the one before it, their per-call times add to an irreducible floor of about 7.7 s. The vision value is measured; the text-summary value is estimated.}
\label{interlude-latency-floor}
\end{figure}

For the representative 37-image workload, local concurrency of six required seven waves, yielding a predicted latency of about 33 seconds. Modeled cloud concurrency of 64 could execute that frontier in one wave; including estimated cold-start and round-trip costs gave about 11.5 seconds, a roughly 2.9× improvement. The prediction depended on those overhead estimates and required staging measurements before cutover. The local cap also reflected memory: concurrent ChunkWorkers each held format-native document state. Skeletonization reduced that per-worker residency, making wider execution feasible without multiplying the same large media working set.

\subsection{Predicting the cost of a design change}\label{heading:predicting-the-cost-of-a-design-change}

A proposed intermediate representation can be inserted into the dependency graph before it is implemented. Recomputing the critical path, execution rounds, peak width, and overhead shows where it could add cost. An IR off the critical path may add no serial latency; one that deepens the chain adds its call time to the floor and may require another distributed round. These are predictions to test, but they help decide which candidate is worth building.

\subsection{Use prediction error as evidence}\label{heading:use-prediction-error-as-evidence}

Implementation eventually supplies measurements, and the useful comparison is not simply whether the result was fast. A residual --- the gap between prediction and observation --- turns into an engineering question, and it drives a loop that runs both before and after implementation (\autoref{interlude-prediction-residual}).

\begin{figure}[htbp]
\centering
\includegraphics[width=\linewidth]{book/latex/assets/interlude-prediction-residual-358697131d.pdf}
\caption{\emph{Using prediction error as evidence.} (a) The gap between a model\textquotesingle s prediction and the observed measurement is a residual, pointing to a wrong parameter, a wrong relation, or a missing mechanism; the numbers are illustrative. (b) Prediction and measurement create a feedback loop: explain the residual, refine the model, and test the revised model against subsequent observations.}
\label{interlude-prediction-residual}
\end{figure}

The queue overhead may have been underestimated. Two apparently independent nodes may share a serialized resource, or media loading may introduce a cost the latency model omitted. A close prediction supplies evidence that the model captures dominant mechanisms; a miss directs investigation toward a wrong parameter, a missing relation, or an implementation problem.

\section{What Carries Forward}\label{heading:what-carries-forward}

DocAble did not turn observed wall-clock time into the durable control. Wall time varies with hosts, services, cold starts, and other environmental effects, so a benchmark score from one machine makes a poor obligation. The performance ratchet holds structural quantities instead --- the quantities the model says contribute to execution cost: production rounds, GenAI calls, capture invocations, and render invocations. Each is counted deterministically on a fixed fixture. A change may lower these quantities freely; raising one requires a named justification recorded beside the baseline. The environment therefore preserves part of the \emph{explanation} for performance rather than a number that only reproduces on the machine that produced it.

A fifth quantity, the number of GenAI hops, was observed without blocking while its behavior was still being characterized. The distinction is worth keeping: a counter that gates and a counter that only reports are different mechanisms, and a new one usually begins as the second.

\subsection{Models get falsified; the obligation survives}\label{heading:models-get-falsified-the-obligation-survives}

A ratchet does not make the current model true forever. Several memory hypotheses in this episode were later falsified. One held that a PDF library\textquotesingle s whole-file object model kept roughly 2.4 GiB resident; measurement found about 1 MiB live, because the file was memory-mapped and the image bytes never entered the managed heap. Another projected that moving rendering off the worker would recover 90 to 120 MiB; the landed measurement recovered about 6, because rendering was already running off the worker. A later adversarial sweep found a scaling dimension the bound had not covered at all: peak memory rose with decoded image area rather than with stored media size. The response was to extend the model and add an admission gate that refuses inputs outside the declared envelope. A capacity model elsewhere in the same environment was eventually caught about twofold too low against dozens of measured sweeps, and its parameter was recalibrated rather than the measurement discounted.

The measurements carried their own provenance. A figure taken on a development host was marked as such and did not become the value of record until the deployment environment reproduced it. A process-level measurement may omit memory included in the container\textquotesingle s total. Recording the measurement\textquotesingle s origin makes that difference visible and helps engineers decide which parameters need to be measured again.

Nor was the memory work finished before the execution-time work began. A memory front was still open while the time models were being built. What became durable was not a particular bound but the requirement to reason explicitly about the resource. Later work could challenge the model, recalibrate it, or replace it with a better one. It could not silently return memory behavior to an unexamined implementation detail. The ratchet preserves the engineering question and the obligation to answer it, even when the answer changes.

% <!-- index-def: ratcheteering -->

% <!-- box-family: canonical -->

\begin{quote}
\mbox{}%
\subsection{Engineering technique --- Ratcheteering}\label{heading:technique-ratcheteering}

Ratcheteering is the practice of converting a consequential engineering concern into explicit models, evidence, and controls that subsequent work must inherit, while keeping those structures open to falsification and revision.

The ratchet preserves the obligation, not necessarily the current answer. A later measurement may refute a model, change a parameter, expose another scaling dimension, or require a different control. The engineering environment should permit those corrections. What later work cannot do silently is return the consequential property to an unexamined implementation detail.

A typical ratchet runs: observe → model → change → test → govern → continue. Later engineering proceeds under the resulting obligation. When new evidence disagrees, revise the model and ratchet again.
\end{quote}

Governance conversion and ratcheteering do different jobs. Governance conversion is the move that turns a recurring failure or judgment into durable structure. Ratcheteering is the practice by which such structures accumulate across episodes while staying correctable.

\section{The Resulting Model Family}\label{heading:the-resulting-model-family}

The large-file requirement did not produce one model. It produced a connected family, each member answering a different question. \autoref{interlude-model-family} collects them.

{\footnotesize
\begin{longtable}[]{@{}
  >{\raggedright\arraybackslash}p{(\linewidth - 2\tabcolsep) * \real{0.5000}}
  >{\raggedright\arraybackslash}p{(\linewidth - 2\tabcolsep) * \real{0.5000}}@{}}
\caption{}\label{interlude-model-family}\tabularnewline
\toprule\noalign{}
\begin{minipage}[b]{\linewidth}\raggedright
Model
\end{minipage} & \begin{minipage}[b]{\linewidth}\raggedright
Question
\end{minipage} \\
\midrule\noalign{}
\endfirsthead
\toprule\noalign{}
\begin{minipage}[b]{\linewidth}\raggedright
Model
\end{minipage} & \begin{minipage}[b]{\linewidth}\raggedright
Question
\end{minipage} \\
\midrule\noalign{}
\endhead
\bottomrule\noalign{}
\endlastfoot
Media residency & What must reside in worker memory? How should memory scale with media size? \\
Batch execution & How do batch size and concurrency affect latency, cost, and truncation risk? \\
Critical path & What sets the latency floor? Which nodes are optimization targets? \\
Distribution & Where should ready work execute? What coordination costs does that introduce? \\
Remediation graph & What work exists, and what depends on what? \\
Fact / IR catalogue & Which derived facts feed later computations? How deep is the information-dependency chain? \\
Image fidelity & What representation of media is required, and what does rendering it cost? \\
\end{longtable}
}

These models take different forms because they answer different questions. Some are typed graphs, some quantitative functions, some declarations of entities and invariants. They stay separate and compose through shared identities and explicit relationships. The remediation graph and fact catalogue supply the static dependency structure, and the batch model supplies per-call cost; the critical-path model combines the two into latency. The distribution model supplies the execution topology the coordination taxes are charged against. The media-residency model supplies the independent memory bound that makes the execution design affordable.

Some answers follow from structure; others depend on calibrated measurements. The critical path identifies ordering that more workers cannot remove, while a latency prediction depends on operation times and coordination costs. Keeping those sources of knowledge distinct makes disagreement informative: a measurement can require a new parameter, a revised relation, or another model.

\subsection{Modeling Inside the Implementation Loop}\label{heading:three-days}

This sequence took about three days inside the existing governed engineering environment. It included skeletonization and streaming, typed-patch work, reconsidering ChunkWorker, remediation and fact graphs, quantitative analysis, the chunkless execution design, and corrections when assumptions proved wrong. The timescale matters because modeling could participate in the implementation loop rather than follow it as a separate documentation phase.

Nobody began with a plan to construct seven models. Larger inputs exposed a memory relationship; changing residency made the execution boundary worth reconsidering; the dependency graph exposed latency limits and possible redesigns. Each representation answered a question the work had made consequential. The earlier examples of successful documents could not by themselves establish a memory bound. That required an account of the scaling relationship, an architectural change, and checks on the assumptions behind it.

Keeping those models inside the loop preserved control while I delegated much of the realization. A black-box optimization process might have reached a faster or smaller implementation sooner, but could have left less evidence about why it worked and what later changes had to preserve. Here, proposed changes could be evaluated against an explicit account and that account could be revised when measurement disagreed. \autoref{interlude-modeling-cycle} shows the continuing interaction among the engineering question, model, analysis, realization, and evidence.

\begin{figure}[htbp]
\centering
\includegraphics[width=\linewidth]{book/latex/assets/interlude-modeling-cycle-9d6747c7c6.pdf}
\caption{\emph{Modeling as a loop.} (a) Models produced after implementation: the linear design → implement → document-with-models sequence. (b) A problem drives a model; analysis yields a prediction and a design; those drive implementation and measurement; measurement refines the model, and the loop continues while the questions justify the cost.}
\label{interlude-modeling-cycle}
\end{figure}

A model need not describe the whole system or survive forever. Keep it while its questions justify its cost. In this episode, the memory obligation survived corrections to the model, and later performance work proceeded within that obligation. This case followed one problem forward as its necessary representations emerged. Brownfield engineering often presents the inverse problem: consequential structure already exists somewhere in the system, but has not been made explicit enough to reason about or govern.
